\documentclass[oneside,a4paper,12pt]{article}
\usepackage[english,brazilian]{babel}
\usepackage[alf]{abntex2cite}
\usepackage[utf8]{inputenc}
\usepackage[T1]{fontenc}
\usepackage[top=20mm, bottom=15mm, left=15mm, right=15mm]{geometry}
\usepackage{framed}
\usepackage{booktabs}
\usepackage{color}
\usepackage{hyperref}
\usepackage{graphicx}
\usepackage{enumerate}
\usepackage{float}
\graphicspath{{./Figuras/}}    
\definecolor{shadecolor}{rgb}{0.8,0.8,0.8}
\pagenumbering{arabic}

\usepackage{multicol}
\setlength{\columnseprule}{1pt}

%Cabeçario
\usepackage{fancyhdr}
\lhead{\rightmark}

\setlength{\parindent}{1.25cm}%paragrafo

\pagestyle{fancy}
\chead{Conteúdo}
\def\MakeUppercase{} %Fonte minúscula no cabeçario

\pagenumbering{arabic}

%Legenda das Figuras
\usepackage[skip=10pt]{caption}
\captionsetup{font={stretch=0.4,small}} 

\newcommand{\frase}{\flushright{Sonhar é acordar-se para dentro.}\\\textbf{Mario Quintana}}

\title{\LINHAHORIZONTAL \\ Título das Nostas de Aula\footnote{Aqui vai o cabeçalho, caso tenha}}
\author{\textbf{Professor Leandro Vieira}\\EREM Regina Pacis\\Palmeirina-PE}
\date{03 de maio de 2020 \\ \LINHAHORIZONTAL}

\newcommand{\LINHAHORIZONTAL}{\rule{17cm}{1.25pt}}

\usepackage{lipsum}

\begin{document}
\maketitle

\tableofcontents

\section{Introdução}

\lipsum[1-2]

\section{Título da Seção}

\lipsum[1-2]

\lipsum[1-2]

\subsection{Novo Sub}

\lipsum[1-2]


	\begin{figure}[h]
	\center
	\includegraphics[width=8cm]{Figuras/imagem.jpg}
	\caption{Legenda da Figura}
	\label{rotulodafigura}
	\end{figure}

\lipsum[1-2]


\section{Título da Seção}

\lipsum[1-2]

\subsection{Novo Sub}

\lipsum[1-2]
	
\subsubsection{Uma Subsubseção} 

\lipsum[1-2]

\end{document}